% -----------------------------------------------------------
\section{Meekness, Meek}
% -----------------------------------------------------------
	\label{MEEKNESS}
	
% ----------------------
\subsection{See also}
% ----------------------

	\hyperref[HUMILITY]{Humility}, \hyperref[LOWLINESS]{Lowliness},

% ----------------------
\subsection{1828 American Dictionary of the English Language} % *1828
% ----------------------
	\label{MEEKNESS-1828}
	
	% -----
	\subsubsection{Meek}
	% -----
	
		\begin{compactitem}
			\item[1] Mild of temper; soft; gentle; not easily provoked or irritated; yielding; given to forbearance under injuries.
			\item[2] Appropriately, humble, in an evangelical sense; submissive to the divine will; not proud, self-sufficient or refractory; not peevish and apt to complain of divine dispensations. Christ says, `Learn of me, for I am meek and lowly in heart, and ye shall find rest to your souls.' Matthew 11:29.
		\end{compactitem}
	
	% -----
	\subsubsection{Meekness}
	% -----

		\begin{compactitem}
			\item Softness of temper; mildness; gentleness; forbearance under injuries and provocations.
			\item[1] In an evangelical sense, humility; resignation; submission to the divine will, without murmuring or peevishness; opposed to pride, arrogance and refractoriness. \hyperref[GALATIANS-5-23]{Galatians 5:23}.
			\item I beseech you by the meekness of Christ. \hyperref[1CORINTHIANS-10-1]{1 Corinthians 10:1}.
			\item Meekness is a grace which Jesus alone inculcated, and which no ancient philosopher seems to have understood or recommended.
		\end{compactitem}

% ----------------------
\subsection{Oxford English Dictionary} % *OED
% ----------------------
	\label{MEEKNESS-OED}
	
	% -----
	\subsubsection{Meek}
	% -----
	
		The ``A'' definitions are for ``meek'' as an adjective; the ``B'' definition is for the word as a noun.

		\begin{compactitem}
			\item[A.1] Gentle, courteous, kind. Of a social superior: merciful, compassionate, indulgent. \textit{Obsolete.}
				\begin{compactitem}
					\item The last entry given here is one from Shakespeare in 1616, so this is an old use, but one that helps us understand its roots.
				\end{compactitem}
			\item[A.2] Humble, submissive; frequently collocated with \textit{mild}.
			\item[A.2.a] Not proud or self-willed; piously humble; patient and unresentful under injury or reproach; (esp. of a woman) demure, quiet. Frequently connoting the gentleness (esp. towards the weak, the humble, and the poor) consonant with a Christian virtue; cf. post-classical Latin \textit{mansuetus} (Vulgate), Hellenistic Greek \gr{πρᾶος} (New Testament). Use of meek as a noun is predominantly in this sense: see branch B. The distinction between this and sense A. 2b is not always clear, since more or less similar qualities of meekness have often been regarded as laudable in certain contexts (as in Christian humility, or, formerly, in women's or servants' characters) but as weak or effeminate in others.
			\item[A.2.b] Humble, submissive (occasionally with \textit{to}). Esp. (\textit{depreciative}): inclined to submit tamely to oppression or injury, easily imposed upon or cowed, timid, biddable.
			\item[A.2.c] In extended use, esp. of flowers: not obtrusively conspicuous, unassuming, modest.
			\item[A.3] Of an animal: tame, domesticated, docile; not fierce. (In early use esp. common in \textit{Scottish})
			\item[A.4] Of an object, substance, the weather, etc.: weak, lacking strength or violence; mild, gentle, soft, light; †supple, pliant (\textit{obsolete}).
			\item[A.5] Proverbial phrases (chiefly in senses A. 1 and A. 2), as \textbf{\textit{as meek as a lamb} (also \textit{a maid}, etc.), \textit{as meek as Moses}}, etc.
			\item[B] With \textit{plural} agreement. Now usually with \textit{the}: meek people as a class. Frequently with allusion (sometimes \textit{ironic}) to Matthew 5:5.
		\end{compactitem}
	
% % ----------------------
% \subsection{Restoration Lexicon} % *RL *RESTORATION LEXICON
% % ----------------------
	% \label{MEEKNESS-OED}

	% \begin{compactitem}
		% \item[]
	% \end{compactitem}

% ----------------------
\subsection{Scriptural Usage} % *SCRIPTURES
% ----------------------
	\label{MEEKNESS-SCRIPTURES}

	% % -----
	% \subsubsection{Old Testament} % *OT
	% % -----
		% \label{MEEKNESS-OT}
	
		% % --
		% \paragraph{KJV}
		% % --
			% \label{MEEKNESS-OT-KJV}
	
			% \begin{tabular}{ll}
				% Total occurrences in the OT & 
			% \end{tabular}
			
			% \begin{compactdesc}
				% \item[]
			% \end{compactdesc}
			
		% % --
		% \paragraph{Hebrew Bible}
		% % --
			% \label{MEEKNESS-HEBREW}
		
			% %
			% \subparagraph{\texorpdfstring{\he{sample word}}{}} % paste actual hebrew text here
			% %
				% \label{PAGA} % whatever is in step bible without periods
			
				% \begin{compactdesc}
					% \item[HALOT , PDF ] \textbf{} % Use the 1994 PDF
						% \begin{compactitem}
							% \item[1]
						% \end{compactitem}
					% \item[BDB , PDF ] \textbf{}
						% \begin{compactitem}
							% \item[1]
						% \end{compactitem}
					% \item[DCH PDF ] \textbf{} % don't include the actual page number, they repeat from volume to volume
						% \begin{compactitem}
							% \item[1]
						% \end{compactitem}
				% \end{compactdesc}
				
				% \begin{compactdesc}
					% \item[Some OT uses] \textbf{}
						% \begin{compactdesc}
							% \item[]
						% \end{compactdesc}
				% \end{compactdesc}
			
		% % --
		% \paragraph{Septuagint}
		% % --
			% \label{MEEKNESS-LXX}
		
			% %
			% \subparagraph{\texorpdfstring{\gr{sample word}}{}}
			% %
		
	% -----
	\subsubsection{New Testament} % *NT
	% -----
		\label{MEEKNESS-NT}
	
		% % --
		% \paragraph{KJV}
		% % --
			% \label{MEEKNESS-NT-KJV}		
	
			% \begin{tabular}{ll}
				% Total occurrences in the NT & 
			% \end{tabular}
			
			% \begin{compactdesc}
				% \item[]
			% \end{compactdesc}
			
		% --
		\paragraph{Greek New Testament} % *GREEK
		% --
			\label{MEEKNESS-GREEK}
		
			%
			\subparagraph{\texorpdfstring{\gr{πραΰς}}{}}
			%
				\label{PRAUS}
			
				\begin{compactdesc}
					\item[BDAG 764a] \textbf{}
						\begin{compactitem}
							\item \textbf{pertaining to not being overly impressed by a sense of one's self-importance, \textit{gentle, humble, considerate, meek}} in the older favorable sense (cp. OED s.v. 1b; Pind., P. 3, 71 describes the ruler of Syracuse as one who is \gr{π}. to his citizens, apparently the rank and file [Gildersleeve]), \textit{unassuming}
								\begin{compactitem}
									\item ``OED 1b'' here seems to refer to OED \#A.1 or \#A.2.a above, I'm not sure.
								\end{compactitem}
						\end{compactitem}
					% \item[CGL , PDF ] \textbf{}
						% \begin{compactitem}
							% \item 
						% \end{compactitem}
					\item[LSJ , PDF ] \textbf{}
						\begin{compactitem}
							\item[I] \textit{mild, soft, gentle}
							\item[I.1] of things
							\item[I.2] of persons, \textit{mild, gentle, meek}
							\item[I.3] of actions, feelings, etc., \textit{mild}
							\item[II] \textit{making mild, taming}
							\item[III] Adverb \gr{πράως}, \textit{mildly, gently}
						\end{compactitem}
						
					\item[TDNT] ...The adv. \gr{πραθως} is often used for the quiet and friendly composure which does not become embittered or angry at what is unpleasant, whether in the form of people (Epict. Ench., 42) or fate. This is an active attitude and deliberate acceptance, not just a passive submission, cf. Epict. Diss., III, 10, 6: \gr{ἂν ἔτι ἐγὼ παρασκευάσωµαι πρὸς τὸ πράως φέρειν τὰ συµβαίνοντα, ὃ θέλει γινέσθω} IV, 7, 12...
						\begin{compactitem}
							\item \textit{Enchiridion} 42: ``Whenever anyone does you wrong or speaks ill of you, remember that he acts and speaks as he does because he thinks it's appropriate for him. He can only conform to his own views, not to yours. So if his views are wrong, he's the one who's harmed, because he's also been deceived. If someone takes a true conjunctive statement to be false, it's not the conjunctive statement that has been harmed but the person who's mistaken. If your inclinations to act are based on these principles, you'll be gentler with anyone who maligns you [\gr{ὁρμώμενος πρᾴως ἕξεις πρὸς τὸν λοιδοροῦντα}], because whenever that happens you'll tell yourself: `That's what he thought it best to do.'''
							\item \textit{Discourses} III.10.6: ``What is it to do philosophy? Isn't it to prepare oneself for whatever happens? Don't you see, then, that what you're saying amounts to: `It's out of my hands whether or not I ever again prepare myself to calmly accept whatever happens [\gr{πρὸς τὸ πρᾴως φέρειν τὰ συμβαίνοντα}]'? This is no different from someone giving up the pancratium because he keeps being hit!''
							\item \textit{Discourses} IV.7.12: ``If someone understands all this, is there anything that could stop him living in a carefree and relaxed way, calmly consenting to everything that might happen to him [\gr{πάντα τὰ συμβαίνειν δυνάμενα πρᾴως ἐκδεχόμενον}] and accepting what has already happened?''
						\end{compactitem}

				\end{compactdesc}
				
				% \begin{tabular}{ll}
					% Total occurrences in the NT & 
				% \end{tabular}
				
				% \begin{compactdesc}
					% \item[Some NT uses] \textbf{}
						% \begin{compactdesc}
							% \item[]
						% \end{compactdesc}
				% \end{compactdesc}
				
			% %
			% \subparagraph{TDNT: \texorpdfstring{\gr{sample word}}{}  (3:536--556)}
			% %
				% \label{KATALLAGE-TDNT}
		
	% % -----
	% \subsubsection{Book of Mormon} % *BOM
	% % -----
		% \label{MEEKNESS-BOM}
	
		% \begin{tabular}{ll}
			% Total occurrences in the BOM & 
		% \end{tabular}
		
		% \begin{compactdesc}
			% \item[]
		% \end{compactdesc}
		
	% -----
	\subsubsection{Doctrine and Covenants} % *DC
	% -----
		\label{MEEKNESS-DC}
	
		% \begin{tabular}{ll}
			% Total occurrences in the DC & 
		% \end{tabular}
		
		\begin{compactdesc}
			\item[{\hyperref[DC97-2]{DC 97:2}}]
		\end{compactdesc}
		
	% % -----
	% \subsubsection{Pearl of Great Price} % *PGP
	% % -----
		% \label{MEEKNESS-PGP}
	
		% \begin{tabular}{ll}
			% Total occurrences in the PGP & 
		% \end{tabular}
		
		% \begin{compactdesc}
			% \item[]
		% \end{compactdesc}
		
% ----------------------
\subsection{Key Phrases} % *KEYPHRASES *PHRASES
% ----------------------
	\label{MEEKNESS-PHRASES}

	\begin{compactdesc}
		\item[meek and lowly] \textbf{12} | \hyperref[MATTHEW-11-29]{Matthew 11:29}; \hyperref[EPHESIANS-4-2]{Ephesians 4:2}; \hyperref[ALMA-37-33]{Alma 37:33}; \hyperref[ALMA-37-34]{Alma 37:34}; \hyperref[MORONI-7-43]{Moroni 7:43}; \hyperref[MORONI-7-44]{Moroni 7:44} (2); \hyperref[MORONI-8-26]{Moroni 8:26} (2); \hyperref[DC32-1]{DC 32:1}; \hyperref[DC107-30]{DC 107:30}; \hyperref[DC118-3]{DC 118:3}
			\begin{compactdesc}
				\item[Bible anchor verses] \phantom{.}
					\begin{compactdesc}
						\item[Matthew 11:29] \gr{ἄρατε τὸν ζυγόν μου ἐφ᾽ ὑμᾶς καὶ μάθετε ἀπ᾽ ἐμοῦ, ὅτι πραΰς εἰμι καὶ ταπεινὸς τῇ καρδίᾳ, καὶ εὑρήσετε ἀνάπαυσιν ταῖς ψυχαῖς ὑμῶν·}
						\item[Ephesians 4:2] \gr{μετὰ πάσης ταπεινοφροσύνης καὶ πραΰτητος, μετὰ μακροθυμίας, ἀνεχόμενοι ἀλλήλων ἐν ἀγάπῃ,}
					\end{compactdesc}
				\item[Commentary] ---
			\end{compactdesc}
	\end{compactdesc}
	
% % ----------------------
% \subsection{Bible Dictionaries} % *DICTIONARIES
% % ----------------------
	% \label{MEEKNESS-BD}

% % ----------------------
% \subsection{Restoration Usage} % *RESTORATION
% % ----------------------
	% \label{MEEKNESS-RESTORATION}

	% % -----
	% \subsubsection{Joseph Smith} % *JS
	% % -----
		% \label{MEEKNESS-JS}
	
		% \begin{compactdesc}
			% \item[{\hyperref[KEY]{Date/Source}}]
		% \end{compactdesc}
		
	% % -----
	% \subsubsection{Temple Ordinances}
	% % -----
		% \label{MEEKNESS-TEMPLE}
	
		% \begin{compactdesc}
			% \item[{\hyperref[KEY]{Date/Source}}]
		% \end{compactdesc}
	
	% % -----
	% \subsubsection{Brigham Young} % *BY
	% % -----
		% \label{MEEKNESS-BY}
	
		% \begin{compactdesc}
			% \item[{\hyperref[KEY]{Date/Source}}]
		% \end{compactdesc}
	
% % ----------------------
% \subsection{Commentary and Studies} % *COMMENTARY *STUDIES
% % ----------------------
	% \label{MEEKNESS-COMMENTARY}

	% % -----
	% \subsubsection{Key Points}
	% % -----
		% \label{MEEKNESS-KEYS}
	
		% \begin{compactitem}
			% \item
		% \end{compactitem}
		
	% % -----
	% \subsubsection{Sample Study}
	% % -----
		% \label{MEEKNESS-study}